%%=====================================================================
%% M2 — WP0 definitions sheet. SURGICAL FRAGMENT (insertion only; no
%% preamble; symbols per mef_macros.tex and the M1 preamble macros
%% \Hquat, \Ad, \rigR, \rigD, \rigM, \rigG). UK English. "We" throughout.
%% Architect draft, 8 Nov 2026 — every choice that is the architect's
%% rather than a PI ruling or an M1 result is marked in a %% ARCH comment.
%% Intended location: M2 Section 2 (after the recall of A1–A4 and M1).
%%=====================================================================

\section{The variables}\label{sec:variables}

The geometry is fixed by M1: $M_{12}=(S^{1}\times\Sigma)\times\Keight$ with
$\Sigma=S^{3}$ and $\Keight=[(\CPtwo\times\Stwo)\times_{w}(\Ttwo/\Ztwo)]^{\Spinc}$,
the fibre $\Ttwo$ carrying the involution of A2 (the fold on the pillowcase; refined in Section~\ref{sec:handedness}) and the
Fibonacci drive $M_{0}$ of A4, with corners $P_{1},\dots,P_{4}$ labelled as in
M1 Appendix~B. We keep every convention of M1 unchanged; in particular the
fibre coordinates are $(\theta_{\psi},\theta_{\omega})$ with half-periods
$(L_{\psi},L_{\omega})$. In this paper $\psi$ and $\omega$ also name
\emph{fields}. The two uses are distinguished by position: a subscript
$\theta_{\psi},\theta_{\omega},L_{\psi},L_{\omega},k_{\omega}$ is a coordinate
label inherited from M1; a bare $\psi$, $\bar\psi$ or $\omega$ is a field
defined in this section. The cumulative table of symbols is Appendix~B.

\subsection{The fifth reading}\label{sec:readings}

M1 numbered four readings by which the four sentences became conditions,
and its abstract already describes the geometry as compact and Riemannian
without making that a numbered reading. This paper makes it one, and
attaches it to every result that depends on it.

\begin{definition}[the metric signs no direction]\label{def:riemannian}
The metric of $M_{12}$ is positive-definite. Time is not a direction the
metric distinguishes by a sign; it is the orbit parameter $s$ of the flow of
A4 on $T\times F$. A derivative ``along time'' is a derivative along that
flow, and every statement below that reads as a statement about time is a
statement along the flow.
\end{definition}

The reading is the geometric form of a sentence as plain as M1's four: the
universe does not label its directions. It is also the reading under which
A1 is unproblematic, since a compact closed manifold with a signed direction
would contain closed time-like curves, and under which the spectra of the
operators of the programme are real and discrete. What singles out a
direction is the drive, and Section~\ref{sec:handedness} shows that the drive
singles it out with an orientation. \rigD
%% PI ruling Q8 (8 Nov 2026). Register: fifth numbered reading of the series.

\paragraph{The standing principle.}
M1 chose among geometries by a minimality principle (M1 Definition~3.2),
and every paper of the series applies the same principle to every datum
the axioms leave open: terms in an action, connections, couplings,
readings. Where a choice exists we take the minimal option and tag the
result \rigD\ with minimality as its hypothesis, on the same footing as
M1's own minimality results; and where the minimal option fails --- a
contradiction, an empty solution set, a conflict with a result already
established --- we say so where it happens rather than enlarge the choice
silently. The present paper applies the principle three times
(Section~\ref{sec:action}), and names each instance.
%% PI, 8 Nov 2026: the one unwritten axiom of the series, now written once.

\subsection{One complex field from two sheets}\label{sec:field}

Axiom A2 supplies a growth law $w=w_{0}e^{\kappa s}$ along the warped
direction and an involution $\iota$ with $\iota^{*}\kappa=-\kappa$
(M1 Definition~2.2). The two signs of $\kappa$ are the two \emph{sheets} of
the exponential, growth and decay, exchanged by $\iota$. M1 identifies
$\iota$ on the fibre algebra: at the drive-fixed corner $P_{1}=C_{0}$ it is
conjugation on $\Hquat$, whose commutative fixed subalgebra containing the
selected unit is $\mathbb{C}_{i}$ (M1 Propositions~5.4 and~5.6); at the three
cycled corners $P_{3}=C_{i}$, $P_{2}=C_{j}$, $P_{4}=C_{k}$ it is the inner
automorphism by the surviving unit, and the drive's three-cycle
$\Ad_{w}:i\to j\to k$ carries one corner's reading to the next
(M1 Proposition~5.6, certificate \texttt{L2\_Ad\_w\_cycles}).

Before the definition, three facts it rests on, each from M1, stated in
words. First, the third axiom says the fibre algebra --- the quaternions,
by M1 Proposition~5.2 --- acts on spinors (M1 Definition~2.5); so whatever
field the two sheets of the growth law make, it is a spinor-type object,
and its values live in the space on which the quaternions act. Second,
that space is fixed by the dimension: in twelve dimensions the Clifford
algebra is the algebra of $32\times32$ quaternionic matrices, by the
eightfold periodicity of Bott (twelve is four more than eight, and the
four-dimensional case is $\Hquat(2)$, so the twelve-dimensional case is
$\Hquat(2)$ enlarged sixteenfold); the space it acts on irreducibly is
therefore thirty-two quaternions, which is sixty-four complex numbers or
one hundred and twenty-eight real ones (\texttt{L7\_module\_dims}). Third,
the operation that exchanges the two sheets is the flow reversal $\tau$ of
Section~\ref{sec:handedness}: reverse the drive and turn the fibre by a
quarter. Done twice it is the fold of the torus that makes the pillowcase,
so on the pillowcase itself it is an involution.

\begin{definition}[the field $\psi$ and its shadow]\label{def:psi}
Let $M\cong\Hquat^{32}$ be the irreducible module of the Clifford algebra
of $M_{12}$. A \emph{sheet field} is a smooth section $\psi$ of $M$: at
every point of the shape, thirty-two quaternions
$\psi_{1},\dots,\psi_{32}$, varying smoothly from point to point. Let
$\hat\tau$ be the lift of the flow reversal $\tau$ to $M$ --- the operation
that carries a spinor at a point to a spinor at the reversed point, and
which, because $\tau$ reverses orientation, is an odd element of the Pin
group. The \emph{shadow} of $\psi$ is the sheet field
\begin{equation}\label{eq:equivariance}
  \bar\psi\;\equiv\;\hat\tau\psi\,,
\end{equation}
the field as seen after the exchange of growth and decay. Applying
$\hat\tau$ twice returns $\psi$ carried round the fold, which on the
pillowcase is $\psi$ itself up to the sign $\pm1$ that M1 Corollary~7.4
assigns to each corner. So on the pillowcase every sheet field splits into
an even part and an odd part under the exchange,
\[
  \psi\;=\;\tfrac12(\psi+\bar\psi)\;+\;\tfrac12(\psi-\bar\psi)\,,
\]
the first unchanged by exchanging growth and decay, the second reversed
by it. Finally, at a corner $C_{u}$ of the pillowcase the quaternions
collapse to the complex numbers $\mathbb{C}_{u}$ fixed by the corner
reading (M1 Propositions~5.4 and~5.6); the \emph{corner reading} of a
sheet field is its image there, thirty-two complex numbers instead of
thirty-two quaternions, and it is in that reading that the phase of
Section~\ref{sec:omega-shape} is defined.
\end{definition}
%% PI rulings 8 Nov 2026 (R-F1(Φ) adopted; Q17(b)): ψ is a section of the
%% Clifford module M ≅ ℍ³² (A3), the bar is the geometric lift τ̂ of the flow
%% reversal WITHOUT conjugation (so ω = ψ + τ̂ψ is quaternion-valued, carrying
%% the imaginary direction along the selected unit — R-τ11), and the
%% corner-algebra bundle 𝒞 is the corner reading of the module-valued field.

The two sheets of A2 are therefore one field: its $\hat\tau$-even part is
what growth and decay share, its $\hat\tau$-odd part is what distinguishes
them. That the field is module-valued rather than real or complex is not an
extra choice: M1 Proposition~7.3 reads the spinor data of A3 as the
$\Spinc$ module, and the selected imaginary unit of Proposition~7.2 is the
direction in $\mathrm{Im}\,\Hquat$ along which the even part carries its
imaginary component. In the corner reading each component is a complex
number $x+iy$ in $\mathbb{C}_{u}$ and $\hat\tau$ sends $y\mapsto-y$.
\rigD

\subsection{The even and odd variables, and the dilaton}\label{sec:omega}

\begin{definition}[$\omega$ and the odd variable]\label{def:omega}
\begin{equation}\label{eq:omega-def}
  \omega\;\equiv\;\psi+\bar\psi\,,
  \qquad
  \psi-\bar\psi\,,
\end{equation}
both sections of $M$. The first is the even part of the sheet field doubled,
the second the odd part doubled. $\omega$ is the Coherent Field of the
programme: it measures how much the two sheets agree. It is quaternion-valued,
and its imaginary part lies along the selected unit; since the odd variable
will enter the dilaton only through its modulus, a non-negative real number,
$\omega$ is the only one of the two that can give the dilaton a direction
(Section~\ref{sec:omega}). In the corner reading, where each component is a
complex number, $\omega=2\,\mathrm{Re}\,\psi$ and $\psi-\bar\psi=2i\,\mathrm{Im}\,\psi$.
\end{definition}

Two properties of $\omega$ are fixed at the level of the definition and are
the inputs to the next proposition. First, $\omega$ is the even
\emph{combination} of the sheets and not their modulus: when growth stands
exactly against decay, $\psi=-\bar\psi$, nothing in the geometry changes and
$\omega=0$. The modulus $|\psi|$ does not vanish there; the combination
does. Second, when the two sheets coincide, $\psi=\bar\psi$, there is no held
distinction between them, and we require that whatever couples the sheets to
the curvature vanish: this is the motivated condition that the coupling is a
coupling \emph{between} the sheets, not a coupling of either sheet alone.
\rigM
%% PI ruling R1 (8 Nov 2026): ω = ψ + ψ̄, with the reason as stated.
%% The second condition reproduces omega-definition-usage_v2.md §4 in-document.

\begin{proposition}[the lowest even invariant]\label{prop:parity}
Write $E=\tfrac12(\psi+\bar\psi)$ for the even part and $O=\tfrac12(\psi-\bar\psi)$
for the odd part of a sheet field. Consider scalars built by multiplying
powers of $E$ and powers of $\lvert O\rvert$ together, and ask for those
that are unchanged when growth and decay are exchanged, and that vanish
both when the sheets coincide ($O=0$) and when growth exactly opposes
decay ($E=0$). Every such scalar has total degree at least three, and the
one of degree three is, up to a real constant, $E\,\lvert O\rvert^{2}$.
In the programme's normalisation it is the dilaton
\begin{equation}\label{eq:PhiQC}
  \QCD\;=\;\omega\,\lvert\psi-\bar\psi\rvert^{2}\,,
\end{equation}
taken component by component in $M$: the $\alpha$-th component of $\QCD$
is $\omega_{\alpha}$ times the squared modulus of $(\psi-\bar\psi)_{\alpha}$.
\end{proposition}

\begin{proof}
We work in the corner reading, where each component is a complex number
$x+iy$ with $x=\mathrm{Re}\,\psi$ the even part and $y=\mathrm{Im}\,\psi$ the
odd part; the general case is the same argument in each component. A
monomial $x^{a}y^{b}$ changes sign under the exchange exactly when $b$ is
odd, because the exchange sends $y$ to $-y$ and leaves $x$ alone; so the
even monomials are those with $b$ even. The monomial vanishes identically
on the locus $y=0$ only if $y$ actually appears, that is $b\ge1$, and on the
locus $x=0$ only if $x$ appears, $a\ge1$. Put the three conditions together:
$b$ is even and at least one, so $b\ge2$; $a\ge1$; hence $a+b\ge3$, with
equality only when $a=1$ and $b=2$. The unique degree-three candidate is
therefore $xy^{2}$. In the original variables $x=\omega/2$ and
$y^{2}=\lvert\psi-\bar\psi\rvert^{2}/4$, so $xy^{2}=\tfrac18\,\omega\lvert\psi-\bar\psi\rvert^{2}$,
which is \eqref{eq:PhiQC} up to the constant $\tfrac18$. The classification
of exponents, the two vanishing criteria and the parity table of $\omega$,
$\psi-\bar\psi$, $\lvert\psi-\bar\psi\rvert^{2}$ and $\QCD$ are certified in
\texttt{M2\_L1\_parity.lean} (\texttt{L1\_unique}, \texttt{L1\_vanish\_Im},
\texttt{L1\_vanish\_Re}, \texttt{L1\_table\_*}). \rigR
\end{proof}

One remark on the module. The factor $\lvert\psi-\bar\psi\rvert^{2}$ is a
non-negative real number in each component, so it carries no direction in
the quaternions: whatever direction $\QCD$ has, it has because $\omega$ has
it. In particular $\QCD$ has an imaginary part only through $\omega$, and
$\omega$ is the unique carrier of the imaginary direction
(\texttt{L7\_imaginary\_square\_real}, which certifies the supporting fact
that the square of an imaginary quaternion is the negative of its squared
modulus). The real factor also commutes with everything, so no question of
the order of factors arises. The quaternionic structure that the fibre
paper of the series reads as the $Sp(1)$ flux on $\Stwo$ is carried by the
direction of $\omega$; this paper does not use it. \rigR

We record what the proposition does and does not decide. Parity alone does
not choose $\psi+\bar\psi$ over the modulus $\lvert\psi\rvert$ as the even
factor: both are unchanged by the exchange. The choice is the first property above, that $\omega$ vanishes
when growth stands against decay; with the modulus in its place $\QCD$ would
not vanish there. The proposition then shows that once the even variable is
$\omega$ and the two vanishing conditions are imposed, there is no further
freedom at lowest order.
%% Correction logged (handover): an earlier argument treating the fold and
%% conjugation as independent involutions was wrong; they are one involution.
%% PI ruling 9 Nov 2026: the square in Φ is the modulus squared |ψ−ψ̄|²
%% (the Source of Truth wrote (ψ−ψ̄)²; Corpus Errors entry 19).

The smallest case by hand: the drive-fixed corner $P_{1}$ is fixed by $\tau$
(Section~\ref{sec:handedness}), so $\bar\psi(P_{1})=\psi(P_{1})$ up to the
corner sign, so $\mathrm{Im}\,\psi(0)=0$ and
$\QCD(P_{1})=0$: the dilaton vanishes at the real corner, where there is one
sheet and nothing to hold apart.

\subsection{What is varied, and over what}\label{sec:scheme}

The field equation of this paper is obtained by variation and reduction. We
state the scheme once so that each later section can cite it.

\begin{definition}[the variational data]\label{def:scheme}
\begin{enumerate}[label=(\roman*),nosep]
\item \emph{The domain} is $M_{12}$ of M1: compact, connected, without
boundary (A1, M1 Definition~2.1).
\item \emph{The fields varied} are the twelve-dimensional metric $g$, with
the A2 warp $w=e^{A}$ depending on the base coordinates
$(\theta_{\psi},\theta_{\omega})$ (M1 Definition~2.2(i)); the dilaton
$\QCD$ of Section~\ref{sec:omega}, a section of the Clifford module $M$ on
$M_{12}$ with coupling $\zeta$ (Section~\ref{sec:equation} states how its kinetic term arises from the sheets);
the $\Spinc$ spinor of A3 on the bundle with determinant line
$L=\mathcal{O}(3)$ (M1 Proposition~7.3); and the $\Spinc$ connection.
\item \emph{The reduction} integrates over $\Keight$ in two stages: over the
six-dimensional fibre $\CPtwo\times\Stwo$, leaving a two-dimensional system
on the pillowcase base $\Ttwo/\Ztwo$ (Section~\ref{sec:warp}); and over the base, leaving
the four-dimensional equation on $S^{1}\times\Sigma$ (Section~\ref{sec:equation}). The field
equation of this paper is the twelve-dimensional one, on $M_{12}$; the
reduction is how its consequences on the base and on the external block
are read.
\item \emph{The scale} of $M_{12}$ is the declared input of M1
(Theorem~8.1, Proposition~8.2) and is not varied.
\end{enumerate}
\end{definition}
%% ARCH: section numbers are placeholders (§3 = WP3, §4 = WP4); renumber at
%% assembly. ROLLBACK 8 Nov 2026: the earlier sentence identifying Φ_QC with
%% the 4D image of the warp was the handover's error and is removed; Φ_QC is
%% a dilaton on M₁₂ (re-scoped R3.1).

